\section{Introdução} 

\subsection{O que é?} % A subsection can be created just before a set of slides with a common theme to further break down your presentation into chunks

%----------------------------------------------
\begin{frame}
	\frametitle{Primeira vez que o termo foi usado}

	\begin{itemize}
		\item Usamos a palavra parede (\textit{Wall}) para indicar proteção contra intrusos;
		\item A primeira vez que o termo firewall foi utilizado para proteger a estrutura e os habitantes de uma construção do fogo;
		\item O maior firewall de todos os tempos;
	\end{itemize}

\end{frame}
%-----------------------------------------------
\begin{frame}
	\frametitle{Denifição}

	\begin{itemize}
		\item Em computação, um firewall é um sistema de segurança de redes de computadores;
		\item Note que é um SISTEMA (pode ser um software, hardware ou um muro mesmo);
		\item SIM, os três podem ser um sistema de firewall em computação (até o muro)!
	\end{itemize}

\end{frame}

%------------------------------------------------

\newcommand{\ondefire}[3]{
	\begin{frame}
		\frametitle{#1}
		\begin{figure}
			\centering
			\includegraphics[height=.5\textheight]{imagens/#3}
			\caption{#2}
		\end{figure}
	\end{frame}
}

\subsection{Onde posso encontrar firewalls?}

\begin{frame}

	\frametitle{É, onde eu posso encontrar esses firewall?}

	\begin{block}{A pergunta que você deve fazer ao sistema que você desconfia ser um firewall é:}

	"Você faz alguma coisa pra proteger meu sistema de ~coisas~ externas?"

	\end{block}


	\textbf{Se a coisa responder "sim", então ele é um firewall!}

\end{frame}

\begin{frame}

	\frametitle{Tudo que protege seu sistema de ameaças externas!}
	
	Você não acredita, né? Então toma uns exemplos:

	\begin{itemize}
		\item{Dispositivos de firewall}
		\item{Software desses dispositivos de firewall}
		\item{Programas do seu sistema operacional}
		\item{Firmware ou Sistema operacional do seu roteador/switch}
		\item{O muro do local físico do servidor (por essa você não esperava, né?)}
		\item{O Zé, o segurança que cuida do prédio onde está localizado o seu seu servidor}
	\end{itemize}

\end{frame}

\subsection{Software}

%ACL
\ondefire{Software}{Access Control List}{acl.png}

%IP TABLES
\ondefire{Software}{iptables}{iptables.jpg}

%ufw Uncomplicated Firewall
\ondefire{Software}{UFW - Uncomplicated Firewall}{ufw.png}

%do_windows
\ondefire{Software}{Firewall do Windows}{dowindows.jpg}

\subsection{Hardware}

%NETGATE
\ondefire{Hardware}{Netgate}{netgate.jpg}

%DELL
\ondefire{Hardware}{Dell}{dell.png}

%Barracuda
\ondefire{Hardware}{Barracuda}{barracuda.jpg}

%Asa
\ondefire{Hardware}{Asa}{asa.jpg}


%------------------------------------------------

\subsection{Analogia do castelo}

%----------------------------------------------
\begin{frame}
	\frametitle{Analogia do Castelo \footnote{Inspirado em \cite{talal}.}}

	%Imagem do Mário----------------------
	\begin{figure}
		\centering
		\includegraphics[height=.75\textheight]{imagens/castelo-mario.png}
		\caption{Castelo sem firewall}
	\end{figure}

\end{frame}

%----------------------------------------------
\newcommand{\analogia}[5][height=.5\textheight]{
	\begin{frame}
	\begin{block}{#2}
		#3
		\begin{figure}
			\centering
			\includegraphics[#1]{imagens/#4}
			\caption{#5}
		\end{figure}
	\end{block}
	\end{frame}
}
\analogia	{O firewall de borda proteje a rede de ataques externos}
		{O posso do castelo protege o castelo de ataques externos}
		{castelo.jpg}
		{O posso do castelo é o "firewall" do castelo}

\analogia	{O firewall pode permitir a entrada/saída da rede e negar a entrada/saída da rede}
		{A ponte do castelo pode ser abaixada para permitir a entrada, e levantada para negar a entrada}
		{castelo-ponte.jpg}
		{A ponte pode recusar sua entrada, Mário.}

\analogia	{Cada host na rede interna tem seu firewall nos sistemas operacionais}
		{Os cidadãos do império tem suas casas com paredes para uma maior seguraça}
		{castelo-analogia.png}
		{Analogia do castelo \cite{tcpip}}

%\item O firewall da borda da rede seria como se fosse o fosso do castelo;
%\item A entrada na rede seria autorizada e fornecida pela ponte;
%\item Cada host ou Sistema Operacional teria sua própria proteção, na analogia seriam as paredes de cada casa;
